% Generated by os-competition-doc diagram generator
% CC-BY-SA 4.0
\documentclass{standalone}
\usepackage{tikz}
\usepackage{amsmath}
\usetikzlibrary{
    arrows.meta,
    shapes.geometric,
    shapes.multipart,
    positioning,
    fit,
    backgrounds,
    calc,
    decorations.pathreplacing,
    shadows,
}
\usepackage{xcolor}
\usepackage{fontspec}
\setmainfont{Noto Sans}



\begin{document}
% ============================================================
% 分层系统架构图模板 (TikZ)
% 参数: LAYERS (层名列表), COMPONENTS (每层组件列表), DATAFLOWS (数据流)
% 图标: cpu, memory, disk, network, ebpf, kernel, userspace, agent
% ============================================================
% ============================================================
% OS竞赛文档 - 学术风格图标库 (TikZ)
% 风格: IEEE 极简线框风
% 配色: 蓝 #00629B / 橙 #C8501E / 灰 #646464
% 使用: \pic[scale=0.8] at (x,y) {cpu};
% ============================================================

% ---------- 基础颜色 ----------
\definecolor{iconblue}{RGB}{0,98,155}
\definecolor{iconorange}{RGB}{200,80,30}
\definecolor{icongray}{RGB}{100,100,100}
\definecolor{iconlight}{RGB}{220,230,240}
\definecolor{iconbg}{RGB}{248,248,248}

\tikzset{
  icon/.style={draw=iconblue, line width=1pt, fill=white},
  iconfill/.style={draw=iconblue, line width=1pt, fill=iconlight},
  iconorange/.style={draw=iconorange, line width=1pt, fill=white},
  icongray/.style={draw=icongray, line width=1pt, fill=white},
  icontext/.style={font=\scriptsize\sffamily, fill=iconblue},
}

% ============================================================
% 1. CPU 芯片
% ============================================================
\tikzset{
  pics/cpu/.style={
    code={
      \draw[icon] (0,0) rectangle (2,2);
      \draw[icon] (0.3,0.3) rectangle (1.7,1.7);
      \foreach \x in {0,0.7,1.4} {
        \foreach \y in {0,0.7,1.4} {
          \draw[icon] (\x+0.4, \y+0.4) circle (0.08);
        }
      }
      % pins left
      \foreach \y in {0.3,0.9,1.5} {
        \draw[icon] (-0.2,\y) -- (0,\y);
      }
      % pins right
      \foreach \y in {0.3,0.9,1.5} {
        \draw[icon] (2,\y) -- (2.2,\y);
      }
      % pins top
      \foreach \x in {0.3,0.9,1.5} {
        \draw[icon] (\x,2) -- (\x,2.2);
      }
      % pins bottom
      \foreach \x in {0.3,0.9,1.5} {
        \draw[icon] (\x,-0.2) -- (\x,0);
      }
    }
  },
  % small variant
  pics/cpu-small/.style={
    code={
      \begin{scope}[scale=0.6, transform shape]
        \pic{cpu};
      \end{scope}
    }
  },
}

% ============================================================
% 2. 内存条
% ============================================================
\tikzset{
  pics/memory/.style={
    code={
      \draw[icon] (0,0) rectangle (3,1);
      \draw[icon] (0.3,0) -- (0.3,1);
      \foreach \x in {0.5,0.7,...,2.7} {
        \draw[icon] (\x,0.15) -- (\x,0.85);
      }
      % notch
      \draw[icon, fill=white] (1.35,0.3) rectangle (1.65,0.7);
      % contacts bottom
      \foreach \x in {0.1,0.3,...,2.9} {
        \draw[icon] (\x,0) -- (\x,-0.15);
      }
    }
  },
}

% ============================================================
% 3. 磁盘/存储
% ============================================================
\tikzset{
  pics/disk/.style={
    code={
      \draw[icon] (0,0) circle (1);
      \draw[icon] (0,0) circle (0.45);
      \draw[iconfill] (0,0) circle (0.35);
      \draw[icon] (0,0) circle (0.12);
      % read head
      \draw[icon] (0.5,0.7) -- (1.2,1.2);
      \draw[icon] (1.0,1.1) circle (0.15);
    }
  },
}

% ============================================================
% 4. 网络节点
% ============================================================
\tikzset{
  pics/network/.style={
    code={
      % globe
      \draw[icon] (0,0) circle (1);
      \draw[icon] (-1,0) arc[start angle=180, end angle=360, x radius=1, y radius=0.4];
      \draw[icon, dashed] (-1,0) arc[start angle=180, end angle=0, x radius=1, y radius=0.4];
      \draw[icon] (0,1) -- (0,-1);
      \draw[icon] (-1,0) -- (1,0);
      % signal
      \foreach \a in {30,60} {
        \draw[iconorange] ({1.2*cos(\a)}, {1.2*sin(\a)}) arc[start angle=\a, end angle=360-\a, radius=1.2];
      }
    }
  },
}

% ============================================================
% 5. 进程/任务
% ============================================================
\tikzset{
  pics/process/.style={
    code={
      \draw[icon, rounded corners=2pt] (0,0) rectangle (1.8,1.5);
      % window controls
      \fill[iconorange] (0.3,1.15) circle (0.08);
      \fill[iconlight, draw=iconblue, line width=0.5pt] (0.6,1.15) circle (0.08);
      \fill[iconlight, draw=iconblue, line width=0.5pt] (0.9,1.15) circle (0.08);
      % content lines
      \foreach \y in {0.8,0.5,0.2} {
        \draw[icongray, line width=0.5pt] (0.4,\y) -- (1.4,\y);
      }
    }
  },
}

% ============================================================
% 6. 锁/同步
% ============================================================
\tikzset{
  pics/lock/.style={
    code={
      % shackle
      \draw[icon, line width=1.2pt] (0.35,0.7) -- (0.35,1.2) arc[start angle=180, end angle=0, radius=0.35] -- (1.05,0.7);
      % body
      \draw[icon, rounded corners=2pt, fill=iconblue!10] (0.15,0) rectangle (1.25,0.75);
      % keyhole
      \draw[icon, fill=white] (0.7,0.3) circle (0.12);
      \draw[icon] (0.7,0.3) -- (0.7,0.5);
      \draw[icon] (0.58,0.5) -- (0.82,0.5);
    }
  },
}

% ============================================================
% 7. eBPF 探针
% ============================================================
\tikzset{
  pics/ebpf/.style={
    code={
      % hook shape
      \draw[iconorange, line width=1.2pt] (1.4,0) -- (0.6,0) -- (0.6,1.2) -- (0,1.2);
      \draw[iconorange, line width=1.2pt] (0,1.3) -- (0,1.5);
      % kernel ring
      \draw[icon] (0.4,0.4) ellipse (0.4 and 0.3);
      \draw[icon] (0.4,0.4) ellipse (0.15 and 0.1);
      % arrow
      \draw[iconorange, line width=0.8pt, ->] (0.1,0.5) -- (0.25,0.45);
    }
  },
}

% ============================================================
% 8. 内核空间
% ============================================================
\tikzset{
  pics/kernel/.style={
    code={
      \draw[icon, rounded corners=3pt] (0,0) rectangle (2.2,1.6);
      \node[icontext] at (1.1,1.25) {\tiny Kernel};
      % tux-like shapes
      \draw[icon] (0.6,0.5) ellipse (0.25 and 0.3);
      \draw[icon] (1.6,0.5) ellipse (0.25 and 0.3);
      \draw[icon, fill=iconblue] (0.55,0.4) circle (0.04);
      \draw[icon, fill=iconblue] (1.55,0.4) circle (0.04);
      \draw[icon] (0.85,0.8) arc[start angle=180, end angle=360, x radius=0.5, y radius=0.2];
      \draw[icon] (1.1,0.65) -- (1.1,0.9);
    }
  },
}

% ============================================================
% 9. 用户空间
% ============================================================
\tikzset{
  pics/userspace/.style={
    code={
      \draw[icon, rounded corners=3pt, fill=iconlight] (0,0) rectangle (2.2,1.6);
      \node[icontext] at (1.1,1.25) {\tiny User};
      % person icon
      \draw[icon] (1.1,0.7) circle (0.2);
      \draw[icon] (0.95,0.5) -- (0.8,0.1) -- (1.4,0.1) -- (1.25,0.5);
      \draw[icon] (0.85,0.1) arc[start angle=180, end angle=360, x radius=0.25, y radius=0.15];
      \draw[icon] (1.1,0.1) arc[start angle=0, end angle=180, x radius=0.25, y radius=0.15];
    }
  },
}

% ============================================================
% 10. Agent/智能体
% ============================================================
\tikzset{
  pics/agent/.style={
    code={
      % head
      \draw[icon, rounded corners=2pt] (0.4,0.9) rectangle (1.6,1.7);
      % eyes
      \fill[iconblue] (0.7,1.35) circle (0.06);
      \fill[iconblue] (1.3,1.35) circle (0.06);
      % antenna
      \draw[icon] (1.0,1.7) -- (1.0,2.0);
      \fill[iconorange] (1.0,2.0) circle (0.08);
      % body
      \draw[icon, rounded corners=2pt] (0.2,0) rectangle (1.8,0.95);
      % body lines
      \foreach \y in {0.3,0.6} {
        \draw[icongray, line width=0.4pt] (0.5,\y) -- (1.5,\y);
      }
    }
  },
}

% ============================================================
% 11. 告警/异常
% ============================================================
\tikzset{
  pics/alert/.style={
    code={
      % triangle
      \draw[iconorange, line width=1.2pt, fill=iconorange!10]
        (1.0,0) -- (0,1.5) -- (2.0,1.5) -- cycle;
      % !
      \draw[iconorange, line width=1.2pt] (1.0,0.5) -- (1.0,1.0);
      \fill[iconorange] (1.0,1.2) circle (0.08);
    }
  },
}

% ============================================================
% 12. 数据库
% ============================================================
\tikzset{
  pics/database/.style={
    code={
      \draw[icon, fill=iconblue!5] (0,0.3) ellipse (0.8 and 0.3);
      \draw[icon] (0,0.3) -- (0,1.2);
      \draw[icon] (0,1.2) ellipse (0.8 and 0.3);
      \draw[icon] (0,0.3) ellipse (0.8 and 0.3);
      \draw[icon] (0,0.8) arc[start angle=0, end angle=360, x radius=0.8, y radius=0.3];
      \draw[icon] (-0.8,0.3) -- (-0.8,1.2);
      \draw[icon] (0.8,0.3) -- (0.8,1.2);
    }
  },
}

% ============================================================
% 13. 文件/文档
% ============================================================
\tikzset{
  pics/file/.style={
    code={
      \draw[icon, fill=white] (0,0) rectangle (1.4,1.8);
      % folded corner
      \draw[icon, fill=iconbg] (1.4,1.4) -- (1.0,1.4) -- (1.0,1.8) -- cycle;
      \draw[icon] (1.0,1.4) -- (1.0,1.8) -- (1.4,1.4);
      % text lines
      \foreach \y in {1.1,0.8,0.5,0.2} {
        \draw[icongray, line width=0.4pt] (0.3,\y) -- (1.1,\y);
      }
    }
  },
}

% ============================================================
% 14. 容器/Docker
% ============================================================
\tikzset{
  pics/container/.style={
    code={
      \draw[icon, rounded corners=3pt] (0,0) rectangle (2.0,1.8);
      % top bar
      \draw[icon, fill=iconblue!8, rounded corners=3pt] (0,1.3) rectangle (2.0,1.8);
      \node[icontext] at (1.0,1.55) {\tiny Container};
      % whale/docker hint
      \draw[icon] (0.5,0.3) rectangle (1.0,0.9);
      \draw[icon] (1.0,0.3) rectangle (1.5,0.9);
      \foreach \x in {0.65,1.15} {
        \draw[icon] (\x,0.9) -- (\x,1.1);
      }
    }
  },
}

% ============================================================
% 15. 图表/监控
% ============================================================
\tikzset{
  pics/monitor/.style={
    code={
      % screen
      \draw[icon, rounded corners=2pt] (0,0.3) rectangle (1.8,1.8);
      % stand
      \draw[icon] (0.5,0) rectangle (1.3,0.3);
      \draw[icon] (0.7,0) -- (0.2,-0.3);
      \draw[icon] (1.1,0) -- (1.6,-0.3);
      % screen content
      \draw[iconorange, line width=0.8pt]
        (0.3,0.5) -- (0.7,1.0) -- (1.1,0.8) -- (1.5,1.4);
      \fill[iconorange] (0.7,1.0) circle (0.06);
      \fill[iconorange] (1.1,0.8) circle (0.06);
    }
  },
}

% ============================================================
% 16. 齿轮/配置
% ============================================================
\tikzset{
  pics/gear/.style={
    code={
      \draw[icon, fill=iconblue!5] (0,0) circle (0.5);
      \foreach \a in {0,45,...,315} {
        \draw[icon, fill=white] ({\a+8.5}:0.44) -- ({\a}:0.65) -- ({\a-8.5}:0.44) -- cycle;
      }
      \draw[icon, fill=white] (0,0) circle (0.2);
    }
  },
}

% ============================================================
% 17. 箭头/数据流（装饰用）
% ============================================================
\tikzset{
  pics/arrow-right/.style={
    code={
      \draw[{Triangle[width=4pt,length=4pt]}-, icon, line width=1.2pt] (0,0) -- (1.8,0);
    }
  },
  pics/arrow-left/.style={
    code={
      \draw[{Triangle[width=4pt,length=4pt]}-, icon, line width=1.2pt] (1.8,0) -- (0,0);
    }
  },
  pics/arrow-down/.style={
    code={
      \draw[{Triangle[width=4pt,length=4pt]}-, icon, line width=1.2pt] (0,1.8) -- (0,0);
    }
  },
  pics/arrow-up/.style={
    code={
      \draw[{Triangle[width=4pt,length=4pt]}-, icon, line width=1.2pt] (0,0) -- (0,1.8);
    }
  },
  pics/arrow-both/.style={
    code={
      \draw[{Triangle[width=4pt,length=4pt]}-{Triangle[width=4pt,length=4pt]}, icon, line width=1.2pt] (0,0) -- (2,0);
    }
  },
}

% ============================================================
% 样式预设：可复用的样式定义
% ============================================================
\tikzset{
  % 架构图层风格
  layer/.style={
    draw=iconblue,
    line width=1.2pt,
    fill=iconlight,
    rounded corners=4pt,
    minimum width=10cm,
    minimum height=1.5cm,
    font=\sffamily\bfseries
  },
  % 组件框
  component/.style={
    draw=iconblue,
    line width=0.8pt,
    fill=white,
    rounded corners=2pt,
    minimum width=2.5cm,
    minimum height=1cm,
    font=\sffamily\small
  },
  % 数据流箭头
  dataflow/.style={
    ->,
    >=Triangle,
    line width=1pt,
    draw=iconorange,
  },
  % 控制流箭头
  controlflow/.style={
    ->,
    >=Triangle,
    line width=0.8pt,
    draw=icongray,
  },
  % 标签节点
  labelnode/.style={
    font=\sffamily\footnotesize,
    fill=white,
    inner sep=2pt,
  },
  % 判断节点（菱形）
  decision/.style={
    draw=iconorange,
    line width=1pt,
    fill=white,
    diamond,
    minimum width=2cm,
    minimum height=1cm,
    font=\sffamily\footnotesize,
  },
  % 起始/结束节点
  terminal/.style={
    draw=iconblue,
    line width=1pt,
    fill=iconblue!10,
    rounded corners=10pt,
    minimum width=2cm,
    minimum height=0.8cm,
    font=\sffamily\footnotesize\bfseries,
  },
}


\begin{tikzpicture}[node distance=0.5cm]

% ---------- 用户可配置的架构参数 ----------
% 层定义: {名称, y坐标, 图标类型, 颜色深度}
% 组件定义: {名称, 层内x偏移, 尺寸}

% ============================================
% LAYER 5: 可视化与告警层 (top)
% ============================================
\node[layer, fill=iconblue!8, minimum height=2.8cm]
  (L5) at (0,8) {};
\node[icontext, font=\large\bfseries] at (-3.5,9.2) {可视化与\& 告警层};

% 组件
\node[component, fill=iconorange!5, draw=iconorange]
  (monitor) at (-0.5,7.0) {};
\pic[scale=1.2] at (-1.5,7.0) {monitor};
\node[labelnode] at (0.8,7.0) {Grafana 仪表盘};

\node[component, fill=iconorange!5, draw=iconorange]
  (alertmgr) at (4,7.0) {};
\pic[scale=1.0] at (3.0,7.0) {alert};
\node[labelnode] at (5.3,7.0) {告警管理};

\draw[dataflow] (-0.5,6.6) |- ++(0,-0.3) -| (4,6.6);

% ============================================
% LAYER 4: 异常检测引擎层
% ============================================
\node[layer, fill=iconblue!5, minimum height=2.5cm]
  (L4) at (0,4.5) {};
\node[icontext, font=\large\bfseries] at (-3.5,5.7) {异常检测引擎层};

\node[component] (rules) at (-2.5,3.5) {};
\pic[scale=0.8] at (-3.3,3.5) {gear};
\node[labelnode] at (-1.0,3.5) {规则引擎};

\node[component] (ml) at (1.5,3.5) {};
\pic[scale=0.8] at (0.3,3.5) {agent};
\node[labelnode] at (2.8,3.5) {ML 检测模型};

\node[component] (scorer) at (5,3.5) {};
\pic[scale=0.8] at (4.2,3.5) {alert};
\node[labelnode] at (6.2,3.5) {异常评分};

\draw[dataflow] (rules) -- (ml);
\draw[dataflow] (ml) -- (scorer);

% ============================================
% LAYER 3: 数据处理与特征提取层
% ============================================
\node[layer, fill=iconlight, minimum height=2.2cm]
  (L3) at (0,1.5) {};
\node[icontext, font=\large\bfseries] at (-3.5,2.5) {数据处理\& 特征提取层};

\node[component] (preproc) at (-2,0.5) {};
\pic[scale=0.8] at (-3.0,0.5) {file};
\node[labelnode] at (-0.5,0.5) {数据预处理};

\node[component] (feat) at (2.5,0.5) {};
\pic[scale=0.8] at (1.3,0.5) {database};
\node[labelnode] at (4.0,0.5) {特征工程};

\draw[dataflow] (preproc) -- (feat);

% ============================================
% LAYER 2 & 1: eBPF数据采集层 (bottom, double)
% ============================================
\node[layer, fill=iconblue!3, minimum height=3.0cm]
  (L2_1) at (0,-1.5) {};
\node[icontext, font=\large\bfseries] at (-3.5,0.0) {eBPF 数据采集层};

% 内核态
\node[component, fill=iconlight] (kern) at (-2,-2.5) {};
\pic[scale=1.0] at (-3.2,-2.5) {kernel};
\node[labelnode] at (-0.8,-2.5) {内核态探针\\(kprobe/tracepoint)};

% 用户态
\node[component, fill=iconlight] (user) at (2.5,-2.5) {};
\pic[scale=1.0] at (1.3,-2.5) {userspace};
\node[labelnode] at (4.0,-2.5) {用户态探针\\(uprobe)};

\draw[dataflow, dashed] (kern) -- (user);

% ============================================
% 层间数据流
% ============================================
\draw[dataflow, ->, line width=1.4pt]
  (L2_1.north) -- (L3.south) node[midway, right] {};
\draw[dataflow, ->, line width=1.4pt]
  (L3.north) -- (L4.south);
\draw[dataflow, ->, line width=1.4pt]
  (L4.north) -- (L5.south);

% 标注: 数据流向
\node[iconorange, font=\sffamily\scriptsize] at (6.2,3.0) {↑ 数据流};
\draw[dataflow, ->] (6.2,-2.8) -- (6.2,8.2);

% ============================================
% 图例
% ============================================
\node[draw=icongray, line width=0.5pt, fill=white, rounded corners=2pt,
      font=\sffamily\tiny, text width=4cm, align=left]
  at (6.5,-3.5) {
  \textbf{图例}\\[2pt]
  \tikz{\draw[dataflow,->] (0,0) -- (0.5,0);}
  数据流 (eBPF事件)\\[2pt]
  \tikz{\draw[controlflow,->] (0,0) -- (0.5,0);}
  控制流\\[2pt]
  \tikz{\draw[dataflow,dashed,->] (0,0) -- (0.5,0);}
  内核态↔用户态交互
};

\end{tikzpicture}

\end{document}
